% !TEX root = tarasov-qasp-zh.tex
% 章首扉图与格言（sec-04，圆周运动）在 zh/ch-07-zh.tex 末尾——照英文源物理位置（ch-07.tex 尾部有效代码）。
\chapter[怎样处理圆周运动问题？]{怎样处理圆周运动问题？\label{ch-08}}

\begin{dialogue}

\tchr 经验告诉我，有关物体圆周运动的问题和习题，对许多应试者来说都是格外困难的。他们对这类问题的回答里包含着大量典型的错误。为了说明这一点，我们请另一位学生来参加讨论。这位学生不了解我们先前讨论过的内容。我们姑且称他为“\textbf{学生乙}”（原先的那位学生此后称为“\textbf{学生甲}”）。

现在请\textbf{学生乙}\hlred{指出：绕地球沿轨道运行的卫星（satellite，俄语称 sputnik）受到哪些力的作用？}我们约定忽略大气的阻力以及月球、太阳和其他天体的引力。

\stdb 卫星受到两个力：地球的引力和离心力（centrifugal force）。

\tchr 地球的引力我没有异议，\hlred{但我不明白你的离心力是从哪里来的。请解释一下。}

\stdb 要是没有这个力，卫星就不可能保持在轨道上。

\tchr 那它会怎么样呢？

\stdb 那还用说，它会掉到地球上来。

\tchr （转向\textbf{学生甲}）记住我先前跟你说过的！这是一个典型的例子：证明某个力存在，不是以物体的相互作用为根据，而是想走“后门”——从物体运动的性质出发。你看，卫星必须待在轨道上，于是就“必须”引入一个拉住它的力。顺便说一句，\hlblue{倘若这个离心力果真存在，卫星反倒无法保持在轨道上，因为作用在卫星上的各力会互相抵消，它将以匀速度沿直线飞行。}
\pnote{学生乙的论证方式——“卫星没有掉下来，所以必有一个拉住它的力”——与第 2 章学生给抛体画“抛力”是同一类错误：从运动的表现反推一个“力”。检验标准仍然只有一条：说得出施力物体（即相互作用）的力才存在。圆周运动中凭空添加“向心力”“离心力”箭头，正是本章要反复纠正的典型错误。}

\stda 离心力从来都不作用在转动的物体上，它作用在连接物（tie，指绳或其他联结构件）上；作用在转动物体上的是向心力。

\stdb 您是说，卫星只受重力作用？

\tchr 对，只受重力。

\stdb 可它竟然不掉到地球上来？

\tchr 在重力作用下的物体的运动就叫下落。所以，卫星确实在下落。只不过，\hlblue{它的“下落”表现为绕地球的圆周运动，因而可以无限期地继续下去。}我们已经确认过，物体运动的方向与它所受力的方向不一定一致（见\chap{ch-04}）。
\pnote{“受重力作用的运动就叫下落，卫星一直在下落”——这个图像与人教版必修第二册第七章《万有引力与宇宙航行》一致：卫星的轨道运动就是一场持续“落向”地球的自由落体，只是水平方向的速度使它始终落不到地面上。“宇宙航行”一节正是由此推出第一宇宙速度 7.9~km/s 的。}

\stdb 我说地球引力和离心力的时候，根据的是公式
\begin{equation}%
\frac{GmM}{r^{2}} = \frac{mv^{2}}{r}
\label{eq-34}
\end{equation}
其中左边是引力，右边是离心力（$v$ 为卫星的速度）；$m$ 为卫星质量，$M$ 为地球质量，$r$ 为轨道半径，$G$ 为引力常量（gravitational constant）。您是说这个公式不对吗？

\tchr 不，公式完全正确，不正确的是你对公式的解释。你把式~\eqref{eq-34} 看成了两个力之间的平衡方程。实际上，它是牛顿第二运动定律
\begin{equation}%
F = ma
\tag{34a}
\label{eq-34a}
\end{equation}
的表达，其中 $F=\dfrac{GmM}{r^{2}}$，而 $a=\dfrac{v^{2}}{r}$ 是向心加速度（centripetal acceleration）。
\pnote{式~\eqref{eq-34} 不是平衡方程：左边是真实存在的万有引力，右边的 $mv^{2}/r$ 不是力，而是 $ma$，故整个等式是 $F=ma$。人教版必修第二册第七章的“万有引力提供向心力”正是此意。}

\stdb 我承认，按您的解释就可以不用离心力了。\hlred{可是，既然没有离心力，总该至少有个向心力吧。}然而您却没有提到这样一个力。

\tchr 在我们的情形里，\hlblue{向心力就是卫星与地球之间的引力。}我要强调，\hlblue{这说的并不是两个不同的力。}绝对不是，这是同一个力。

\stdb 那为什么还要引入向心力这个概念呢？

\tchr 在这一点上我完全同意你的看法。依我看，“向心力”这个术语只会带来混乱。所谓向心力，根本不是与别的力一起作用在物体上的一个独立的力；\hlblue{它是作用在沿圆周匀速运动的物体上的一切力的合力。}

量 $mv^{2}/r$ 并不是一个力。它等于物体的质量 $m$ 与向心加速度 $v^{2}/r$ 的乘积。这个加速度指向圆心，因此，作用在沿圆周匀速运动的物体上的一切力的合力也指向圆心。可见，\hlblue{存在的是向心加速度，以及合成后使物体获得这个向心加速度的那些力。}
\pnote{与人教版对照：现行教材保留了“向心力”一词（必修第二册第六章 6.2“向心力”），但同样明确它是按效果命名的力——由其他力或它们的合力提供，受力分析时不能作为额外的箭头画上去。教师“与其用‘向心力’，不如直接说合力”的主张是原书的教学立场，与教材精神一致：两种说法下，$mv^{2}/r$ 都不是独立的力。}

\stdb 我得承认，这种处理圆周运动的观点我很赞成。这种运动确实不是以力的平衡为特征的静力学问题，而是动力学问题。

\stda 如果我们摒弃向心力的概念，那么“离心力”这个词恐怕也该一并抛弃，哪怕是在谈到连接物时。

\tchr \hlblue{“离心力”这个术语的引入就更没有道理了。}向心力毕竟还是实际存在的——尽管只是作为合力存在。\hlblue{而离心力在许多情形下根本就不存在。}

\stda 您最后一句话我不明白。离心力是作为向心力的反作用引入的。如果照您所说，它并不总存在，那牛顿第三运动定律不就不是普遍成立的了吗？是这样吗？

\tchr \hlblue{牛顿第三定律只对由物体相互作用决定的真实的力成立，而对它们的合力不成立。}我可以借你们已经熟悉的圆锥摆（conical pendulum）为例来说明（\figr{fig-33}）。

\begin{figure}[!ht]
\centering
\includegraphics[width=0.6\textwidth]{figs/ch-08/fig-033a.pdf}
\caption[圆锥摆。]{圆锥摆。系在绳上的物体在水平面内做圆周运动；运动过程中它划出的立体具有圆锥的形状。\label{fig-33}}
\end{figure}

小球受两个力：重力 $P$ 和绳的张力 $T$。这两个力合在一起使小球获得向心加速度，它们的和就称为向心力。力 $P$ 源于小球与地球的相互作用，它的反作用力是作用在地球上的力 $P_{1}$；力 $T$ 源于小球与绳的相互作用，它的反作用力是作用在绳上的力 $T_{1}$。

如果把 $P_{1}$ 与 $T_{1}$ 形式地相加，得到的就是通常所理解的离心力（见\figr{fig-33} 中的虚线）。可是这个力作用在哪里呢？它的一个分量作用在地球上，另一个却作用在完全不同的物体——绳——上，这时我们还有理由把它叫做一个力吗？显然，在这种情形里，离心力概念没有任何物理意义。

\stda 那么在哪些情形下离心力存在呢？

\tchr 比如沿轨道运行的卫星，这里只有地球和卫星两个物体相互作用。向心力是地球吸引卫星的力，\hlblue{离心力则是卫星吸引地球的力。}
\pnote{本书的“离心力”大体指“向心力的反作用”，且仅见于地球—卫星两体问题（卫星吸引地球的力）；它与大学在转动参考系中引入的“惯性离心力”（无施力物体）是两回事。人教版不用“离心力”一词，只在必修第二册第六章“生活中的圆周运动”讲离心现象——合力突然消失或不足以提供向心力时，物体沿切线飞出或远离圆心。}

\stdb 您说过，\hlblue{牛顿第三定律对真实力的合力不成立。}我想，在这种情形下，它对一个真实力的分力同样也不成立。是这样吗？

\tchr 对，完全正确。在这方面我举一个与转动毫不相干的例子。一个球放在地面上，并与一面同地面成钝角的墙相接触（\figr{fig-34}）。

\begin{figure}[!h]%
\centering
\includegraphics[width=0.4\textwidth]{figs/ch-08/fig-034a.pdf}
\caption[球静置于地面，触及与地面成钝角的墙。]{球静置于地面，并与一面同地面成钝角的墙相接触。\label{fig-34}}
\end{figure}

把球的重力分解为两个分力：一个垂直于墙面，一个平行于地面。我们就用这两个分力来代替球的重力。倘若牛顿第三定律适用于各个分力，那就可以期望墙给球一个反作用力，与垂直于墙的那个重力分力相平衡。于是，平行于地面的重力分力便无人平衡，球就该获得水平加速度了。显然，这在物理上是荒谬的。

\stda 到目前为止，你们讨论的都是匀速圆周运动。\hlred{物体沿圆周做非匀速运动时，该怎么处理？}比如，一个物体从竖直固定的圆环顶端滑下。它沿环滑动时是在做圆周运动，但这不可能是匀速的，因为物体的速度在不断增大。这种情形下该怎么办呢？

\tchr 如果物体沿圆周匀速运动，作用在它上面的所有力的合力必定指向圆心，使物体获得向心加速度。在较一般的非匀速圆周运动情形，合力并不严格指向圆心。这时，它有一个沿半径指向圆心的分量，还有一个沿物体轨迹（即圆周）切线方向的分量。前一个分量负责物体的向心加速度，后一个分量负责所谓的切向加速度（tangential acceleration），它与速度的变化相联系。应当指出，既然物体的速度在变，向心加速度 $v^{2}/r$ 也必然随之改变。
\pnote{人教版必修第二册第六章只处理匀速圆周运动（6.4“向心加速度”）。物体沿圆周变速运动时，合力一般不再指向圆心：切向分量改变速率，径向（法向）分量改变方向——即竞赛与大学普通物理中的“自然坐标系”（切向—法向分解）。式 $a=v^{2}/r$ 每一时刻都成立，只需代入瞬时速率。}

\stda 这是不是说，每一时刻的向心加速度都由公式 $a = v^{2}/r$ 确定，其中 $v$ 是该时刻的瞬时速度？

\tchr 正是如此。匀速圆周运动的向心加速度是恒定的，而非匀速圆周运动的向心加速度在运动过程中不断变化。

\stda 那么，怎样才能知道在非匀速转动中速度 $v$ 究竟如何变化呢？

\tchr 通常为此要借助于能量守恒定律（law of conservation of energy）。我们来看一个具体的例子。
\hlred{设一物体从竖直放置、半径为 $R$ 的圆环顶端无摩擦地滑下。当它经过比环顶低 $h$~\si{\centi\metre} 的一点时，它压在环上的力是多大？物体在环顶时的初速度为零。}

首先，必须弄清有哪些力作用在物体上。
\begin{figure}[!ht]%
\centering
\includegraphics[width=0.6\textwidth]{figs/ch-08/fig-035a.pdf}
\caption[非匀速圆周运动中物体所受各力的分解。]{物体做非匀速圆周运动时所受各力的分解。\label{fig-35}}
\end{figure}
\stda 物体受两个力：重力 $P$ 和支承力 $N$，如\figr{fig-35} 所示。

\tchr 对。接下来你打算怎么做？

\stda 我就按您说的做。先求这两个力的合力，再把合力分解为两个分量：一个沿半径，另一个沿圆周的切线方向。

\tchr 完全正确。不过，更简单的做法显然是一开始就把作用在物体上的两个力沿这两个方向直接分解，而不必先求合力——何况真正需要分解的只有一个力，即重力。

\stda 我对各力的分解如\figr{fig-35} 所示。

\tchr 力 $P_{2}$ 决定物体的切向加速度，眼下我们对它不感兴趣。力 $P_{1}$ 与 $N$ 的合力使物体获得向心加速度，即
\begin{equation}%
P_{1} - N = \frac{mv^{2}}{R}
\label{eq-35}
\end{equation}
物体在我们关心的那一点（\figr{fig-35} 中的点 $A$）的速度，可由能量守恒定律求得：
\begin{equation}
Ph = \frac{mv^{2}}{2}
\label{eq-36}
\end{equation}
把式~\eqref{eq-35} 与式~\eqref{eq-36} 联立，并考虑到
\begin{equation*}%
P_{1}=P \cos \alpha = P \left(\frac{R-h}{R}\right),
\end{equation*}
可得
\begin{equation*}
\frac{P}{R} \left( R-h \right) - N = \left(\frac{2Ph}{R}\right)
\end{equation*}
根据牛顿第三定律，所求的物体压在环上的力等于支承力
\begin{equation}%
N=P \left( \frac{R-3h}{R} \right)
\label{eq-37}
\end{equation}
\stdb 您这里假定了物体到达点 $A$ 时仍在环面上。可是它可能在到达点 $A$ 之前就飞离圆环了。

\tchr 我们可以求出物体离开环面的那个点。这个点对应于物体压环的力减小到零的极限情形。因此，在式~\eqref{eq-37} 中令 $N=0$，解出 $h$，即从环顶到物体飞离之点的竖直距离。于是
\begin{equation}%
h_{0} = \frac{R}{3}
\label{eq-38}
\end{equation}
若题给的 $h$ 满足条件 $h < h_{0}$，则式~\eqref{eq-37} 的结果成立；反之，若 $h\geqslant h_{0}$，则 $N=0$。
\pnote{$h_{0}=R/3$（脱离点与圆心的连线同竖直方向的夹角满足 $\cos\alpha=2/3$）是“物体从光滑球面或环面外侧滑下”这类问题的经典结论，判据就是教师所说的 $N=0$——此后仅由重力的径向分量提供向心加速度。这个模型在高考与竞赛中反复出现（如小球从光滑半球顶滑下求脱离位置）。}

\stda 据我理解，解这道题用到了两条物理定律，即式~\eqref{eq-35} 和式~\eqref{eq-36}。

\tchr 你能指出这一点非常好。确实，这道题的求解涉及两条定律：牛顿第二运动定律（见式~\eqref{eq-35}）和能量守恒定律（见式~\eqref{eq-36}）。遗憾的是，应试者在解题时并不总能清楚地意识到自己用的究竟是哪几条物理定律。而这一点，我认为至关重要。
\pnote{“动力学方程（沿半径方向的牛顿第二定律）＋能量方程”联立，是中学处理变速圆周运动的标准两步法。能量方程的依据——机械能守恒定律——见人教版必修第二册第八章《机械能守恒定律》；本章各题均为无摩擦、只有重力做功的情形，正满足机械能守恒的条件。}

再看下面这个例子。
\hlred{给物体一个初速度 $v_{0}$，使它从 $A$ 点移动到 $C$ 点。现有两条可供选择的路径\drkgry{（\figr{fig-36}~(a)、(b)）}。两种情形中物体都要到达同样的高度 $H$，只是方式不同。求两种情形各自所需的最小初速度 $v_{0}$。摩擦可以忽略。}

\stdb 我想两种情形下的最小初速度应当相同，因为没有摩擦，要到达的高度又相同。这个速度可以由能量守恒定律算出：
\begin{equation*}%
mgH = \frac{mv_{0}^{2}}{2} \, \, \text{由此} \, \, v_{0} = \sqrt{2gH}
\end{equation*}
\begin{figure}[!h]%
\centering
\includegraphics[width=0.42\textwidth]{figs/ch-08/fig-036a.pdf}
\caption[物体沿两条不同路径从 A 点到 C 点。]{物体沿两条不同的路径从点 $A$ 运动到点 $C$。问题是求两种情形下所需的最小初速度。\label{fig-36}}
\end{figure}

\tchr 你的回答错了。你忽略了一个事实：在第一种情形里，物体通过轨迹最高点时，仍处在转动状态之中。这就是说，在最高点 $B$（\figr{fig-36}~\drkgry{(a)}），物体应有一个速度 $v_{1}$，它由与式~\eqref{eq-35} 类似的动力学方程确定。既然题中要求的是最小值，就应考虑物体在点 $B$ 压紧支承物的力减小到零的极限情形。这时作用在物体上的只有重力，由重力来赋予物体向心加速度。于是
\begin{equation}%
mg=\frac{mv_{1}^{2}}{R} = \frac{2mv_{1}^{2}}{H}
\label{eq-39}
\end{equation}
在动力学方程~\eqref{eq-39} 之外再补上能量方程
\begin{equation}%
\frac{mv_{0}^{2}}{2} = \frac{mv_{1}^{2}}{2} + mgH
\label{eq-40}
\end{equation}
便求得最小初速度等于 $\sqrt{5gH/2}$。在第二种情形里，物体可以用无限接近于零的速度通过最高点，因此只用能量方程就够了。这时你的答案是正确的。

\stdb 现在我明白了。第一种情形中，要是物体在点 $B$ 没有速度，它就会直接从轨道上掉下来。

\tchr 假如第一种情形按你的建议取初速度 $v_{0}=\sqrt{2gH}$，物体根本到不了点 $B$，而会在此前某处就脱离轨道。我建议你们\hlred{求出当初速度为 $v_{0}=\sqrt{2gH}$ 时，物体脱离轨道那一点的高度 $h$。}

\stda 请让我来试着解这道题。

\tchr 当然可以。

\begin{figure}[!ht]
\centering
\includegraphics[width=0.5\textwidth]{figs/ch-08/fig-037a.pdf}
\caption[在圆形轨道上运动的物体。]{在圆形轨道上运动的物体。问题是当初速度为 $v_{0}=\sqrt{2gH}$ 时，求物体将脱离轨道的那一点。\label{fig-37}}
\end{figure}
\stda 在物体脱离轨道的那一点，支承力显然等于零，因此该点处物体只受重力作用。可以把重力分解为两个分量：一个沿半径方向（$mg \cos \alpha$），另一个垂直于半径方向（$mg \sin \alpha$），如\figr{fig-37} 所示（点 $A$ 即物体脱离轨道之点）。沿半径的分量使物体获得向心加速度，由如下方程确定：
\begin{equation}%
mg \cos \alpha = \frac{mv_{2}^{2}}{R}
\label{eq-41}
\end{equation}
\pnote{底本式~\eqref{eq-41} 误排为 $\frac{mv_{2}^{2}}{2}$：沿半径方向的牛顿第二定律应为 $\frac{mv_{2}^{2}}{R}$；也唯有如此，与式~\eqref{eq-42} 联立并代入 $\cos \alpha = (h-R)/R$，才能得出下文的 $mg(h-R) = mv_{0}^{2} - 2mgh$。中文版径改。}
其中 $v_{2}$ 是物体在点 $A$ 的速度。为求出它，可以利用能量方程
\begin{equation}%
\frac{mv_{2}^{2}}{2} + mgh = \frac{mv_{0}^{2}}{2}
\label{eq-42}
\end{equation}
把动力学方程~\eqref{eq-41} 与能量方程~\eqref{eq-42} 联立，并考虑到 $\cos \alpha = (h-R)/R$，得
\begin{equation*}%
mg(h-R) = mv_{0}^{2} - 2 mgh
\end{equation*}
由此
\begin{equation}%
h = \frac{2v_{0}^{2}+gH}{6g}
\label{eq-43}
\end{equation}
代入 $v_{0}^{2}=2gH$，最终结果为
\begin{equation*}%
h = \frac{5}{6} H
\end{equation*}

\tchr 完全正确。注意，利用式~\eqref{eq-43} 也可以求出使物体翻环（loop the loop，即翻越整个竖直圆环）的初速度 $v_{0}$：在式~\eqref{eq-43} 中取 $h=H$，于是
\begin{equation*}%
H = \frac{2v_{0}^{2}+gH}{6g}
\end{equation*}
由此直接得到先前已经求出的结果
\begin{equation*}
v_{0} = \sqrt{\frac{5gH}{2}}
\end{equation*}
\pnote{把 $H=2R$（圆环直径，即最高点高度）代入，即得著名的翻环临界条件：最高点速度 $v=\sqrt{gR}$，相当于从高度 $5R/2$ 处由静止滑下。这就是游乐场“过山车”与杂技“水流星”的物理原理——竖直圆周内侧轨道最高点的最小速度问题，在高考与竞赛中极为常见。}

\stda 条件~(\ref{eq-43}) 是就物体脱离轨道的情形得出的。怎么能把它用于物体不脱落而翻过整个圆环的情形呢？

\tchr 恰在环的最高点“脱落”，实际上意味着物体并未脱落，而是通过该点继续沿圆周运动。

\stdb 也可以说，物体好像只在某一瞬间脱落了一下。

\begin{figure}[!h]%
\centering
\includegraphics[width=0.4\textwidth]{figs/ch-08/fig-038a.pdf}
\caption[物体静置于旋转的斜面上。]{物体静置于旋转的斜面上。问题是求使物体不致从斜面滑下所需的静摩擦系数的最小值。\label{fig-38}}
\end{figure}
\pnote{底本此图图题作“求使物体留在斜面上不滑下所需静摩擦系数的最大值（maximum value）”，与正文“求最小系数 $k_{0}$”相矛盾（$k_{0}$ 表征的是静摩擦力的最大可能值），显系转写之误；中文版图题按正文拟改，待考原书。}

\tchr 说得对。最后我提出下面这个问题\drkgry{（\figr{fig-38}~(a)）}。
\hlred{物体位于倾角为 $\alpha$ 的斜面的下端。斜面绕竖直轴以角速度（angular velocity）$\omega$ 匀速转动，物体到转轴的距离为 $R$。求使物体能保持在旋转斜面上而不滑下所需的最小系数 $k_{0}$——我提醒一下，这个系数表征静摩擦力（force of static friction）的最大可能值。}

我们照例从这个问题开始：有哪些力作用在物体上？

\stda 物体受三个力：重力 $P$、支承力 $N$ 和摩擦力 $F_{fr}$。

\tchr 完全正确。好在你没有再加上一个向心力。下一步打算怎么做？

\stda 接下来，我把各力沿斜面方向和垂直斜面方向分解，如\figr{fig-38}~\drkgry{(b)} 所示。

\tchr 请容许我在这里打断你。我不喜欢你这种分解力的办法。请告诉我：物体沿哪个方向加速？

\stda 加速度沿水平方向，是向心加速度。

\tchr 好。正因为如此，才应当把力沿水平方向（即沿加速度方向）和竖直方向（即垂直于加速度方向）分解。还记得我们在\chap{ch-06} 里讨论过的内容吧。

\stda 现在我明白了。各力沿水平方向和竖直方向的分解如\figr{fig-38}~{(c)} 所示。各力的竖直分量相互平衡，水平分量则使物体获得加速度。于是
\begin{equation*}%
\left.
\begin{aligned}
N \cos \alpha + F_{fr} \sin \alpha = P \\
F_{fr} \cos \alpha -N \sin \alpha = \frac{mv^{2}}{R}
\end{aligned}
\right\}
\end{equation*}

考虑到
\begin{equation*}%
F_{fr}=k_{0}N, \, \, \left( \frac{v^{2}}{R} \right)=\omega^{2}R \, \, \text{且} \, \, m=\left(\frac{P}{g} \right),
\end{equation*}
可把这些方程改写为
\begin{equation*}%
\left.
\begin{aligned}
N (\cos \alpha + k_{0} \sin \alpha)=P \\
N(k_{0} \cos \alpha -\sin \alpha)=\frac{P \omega^{2} R}{g}
\end{aligned}
\right\}
\end{equation*}
\pnote{底本此式右端误排为 $P \omega^{2} Rg$（乘 $g$）。由 $m=P/g$ 与 $a=\omega^{2}R$，水平方向方程的右端应为 $P \omega^{2} R/g$（除以 $g$）；这也与下文两式相除后所得的 $g/(\omega^{2}R)$ 相一致。中文版径改。}

\stdb 可是你只有两个方程，却有 $k_{0}$、$P$、$N$ 三个未知量。\\
\tchr 这不碍事。我们并不需要求出全部三个未知量，只要求出系数 $k_{0}$。用第一式除以第二式，未知量 $P$ 和 $N$ 就都容易消去。

\stda 相除后得
\begin{equation*}%
\frac{\cos \alpha+ k_{0} \sin \alpha}{k_{0} \cos \alpha - \sin \alpha} = \frac{g}{\omega^{2}R}
\end{equation*}
解出所求的系数
\begin{equation}%
k_{0} =\left( \frac{\omega^{2}R \cos \alpha+g \sin \alpha}{g \cos \alpha- \omega^{2} R \sin \alpha} \right)
\label{eq-44}
\end{equation}
\tchr 由式~\eqref{eq-44} 显然可见，条件
\begin{equation*}%
\left( g \cos \alpha- \omega^{2}R \sin \alpha \right) >0
\end{equation*}
必须成立。这个条件也可以写成
\begin{equation}%
\tan \alpha < \frac{g}{\omega^{2}R}
\label{eq-45}
\end{equation}
若条件~\eqref{eq-45} 不满足，那么无论多大的摩擦力都无法把物体保持在旋转的斜面上。

\end{dialogue}

\section*{习题}
\label{prob-08}
\addcontentsline{toc}{section}{习题}

% 注：习题系列编号依赖第 5 章的 series=problems 声明（ch-05-zh 尚未译出前，
% 编译此列表会报 enumitem“未定义系列”错误，属暂态，译出 ch-05-zh 后自动消失）。
\begin{enumerate}[resume*=problems]
\item 一辆坦克压在凸形桥桥面中点和压在凹形桥桥面中点的力之比是多少？两种情形下桥的曲率半径（radius of curvature）均为 \SI{40}{\metre}，坦克的速度为 \SI[per-mode=symbol]{45}{\kilo\metre\per\hour}。
\pnote{坦克过凸桥顶与凹桥底（人教版必修第二册 6.5“生活中的圆周运动”的经典情境，教材以汽车过拱形桥为例）：凸桥顶 $N=mg-mv^{2}/r<mg$，凹桥底 $N=mg+mv^{2}/r>mg$。}
\item 物体从 $H=\SI{60}{\centi\metre}$ 的高度无摩擦地滑下，随后沿半径 $R=\SI{20}{\centi\metre}$ 的圆形轨道翻环（\figr{fig-39}）。求物体在 $A$、$B$、$C$ 三点压在轨道上的力之比。
\begin{figure}[!ht]
\centering
\includegraphics[width=0.5\textwidth]{figs/ch-08/fig-039a.pdf}
\caption[物体滑下后沿圆轨道翻环。]{物体从给定高度滑下，随后沿圆形轨道翻环。问题是求物体在 $A$、$B$、$C$ 三点压在轨道上的力之比。见习题 16。\label{fig-39}}
\end{figure}
\item 物体系在长为 $R$ 的绳端，可在竖直平面内转动。应使物体在最高点获得多大的水平速度，才能使它在最低点时绳的张力等于物重的十倍？
\item 一颗卫星绕某球形行星沿圆轨道转动，周期为 $T$，轨道到行星表面的距离等于行星半径 $R$ 的一半。求行星物质的密度。引力常量记为 $G$。
\item 质量为 $m$ 的物体可沿弯成半径为 $R$ 的圆弧形槽无摩擦滑动。当槽绕竖直轴以角速度 $\omega$ 匀速转动时（\figr{fig-40}），物体将在多高 $h$ 处相对槽保持静止？物体对槽的压力 $F$ 是多大？
\begin{figure}[!ht]
\centering
\includegraphics[width=0.35\textwidth]{figs/ch-08/fig-040a.pdf}
\caption[物体在匀速转动的圆弧槽内滑动。]{物体在以匀角速度转动的圆弧形槽内无摩擦滑动。问题是求物体对槽的压力以及物体相对槽静止时所在的高度。见习题 19。\label{fig-40}}
\end{figure}
\item 半径为 $R$ 的圆环竖直固定在地面上。物体从环顶无摩擦地滑下（\figr{fig-41}）。物体将落在离环的固定点多远的距离 $l$ 处？
\begin{figure}[!ht]
\centering
\includegraphics[width=0.35\textwidth]{figs/ch-08/fig-041a.pdf}
\caption[物体从竖直固定的圆环顶端滑下。]{物体从竖直固定于地面的圆环顶端无摩擦滑下。问题是求物体落地点与环的固定点之间的距离 $l$。见习题 20。\label{fig-41}}
\end{figure}
\end{enumerate}

% NEWTERMS: centrifugal force | 离心力 | 本章核心辨析对象；中学不使用该词，只讲“离心现象”
% NEWTERMS: centripetal acceleration | 向心加速度 | 人教版必修第二册 6.4
% NEWTERMS: tangential acceleration | 切向加速度 | 变速圆周运动；自然坐标系（切向—法向）之切向
% NEWTERMS: conical pendulum | 圆锥摆 | 
% NEWTERMS: angular velocity | 角速度 | 
% NEWTERMS: gravitational constant | 引力常量 | 原书符号 $G$；人教版称引力常量
% NEWTERMS: law of conservation of energy | 能量守恒定律 | 人教版必修第二册第八章《机械能守恒定律》
% NEWTERMS: force of static friction | 静摩擦力 | 原书 $k_0$ 表征其最大可能值
% NEWTERMS: coefficient of static friction | 静摩擦系数 | 原书符号 $k_0$；与第 2 章滑动摩擦系数 $k$ 相应
% NEWTERMS: radius of curvature | 曲率半径 | 习题用词
% NEWTERMS: tie (bonding member) | 连接物 | 指绳、杆等联结构件
% NEWTERMS: circular motion | 圆周运动 | 第 2 章已用，glossary 补录
% NEWTERMS: loop the loop | 翻环（翻越竖直圆环）| 习语性表达，非正式术语
